\chapter{2022:Carolyn Bertozzi and Morten Meldal and K. Barry Sharpless}

\label{chap:chemistry2022}

The Nobel Prize in Chemistry for the year 2022 was awarded jointly to Carolyn Bertozzi and Morten Meldal and K. Barry Sharpless.

\textbf{Motivation:}

for the development of click chemistry and bioorthogonal chemistry

\section{Carolyn Bertozzi}

\subsection*{Early Life and Education}

Carolyn Bertozzi was born in 1966 in Boston, Massachusetts, USA.

\subsection*{Career and Major Research}

Bertozzi earned a bachelor's degree from Harvard University in 1988 and a doctorate from the University of California, Berkeley, in 1993. She became a professor at Stanford University and an investigator of the Howard Hughes Medical Institute. She developed bioorthogonal chemistry, reactions that proceed inside living organisms without disturbing the surrounding chemistry, and used them to study the sugars on cell surfaces.

\subsection*{Nobel Prize-winning Work}

The work for which Carolyn Bertozzi was awarded the Nobel Prize in Chemistry in 2022 was described by the Nobel Committee as: "for the development of click chemistry and bioorthogonal chemistry".

Bertozzi's modified Staudinger reaction joined a chemical tag on a cell surface sugar to a probe inside a live organism, demonstrating that bioorthogonal reactions can be carried out in living systems.

\subsection*{Significance and Impact}

Her methods allow researchers to follow biomolecules such as glycans in living cells and animals, and they have been applied to the delivery of drugs and imaging agents.

\subsection*{Later Career and Honors}

Bertozzi continues her research on chemical biology at Stanford University, where she directs the Sarafan ChEM-H institute.

\subsection*{Legacy}

Bioorthogonal chemistry has become a standard tool for labeling and imaging molecules in vivo and is used in the design of medicines.

\subsection*{Modern Developments}

Bertozzi's development of bioorthogonal chemistry and its application to imaging glycans on living cells founded the field of bioorthogonal chemistry and modern chemical biology. Her methods, including click chemistry in living systems, are now used to study glycosylation, develop drug conjugates, and enable in vivo imaging, with applications in medicine and materials.

\subsection*{Selected Publications}

"Cell Surface Engineering by a Modified Staudinger Reaction" (Science, 2000)

\clearpage

\section{Morten Meldal}

\subsection*{Early Life and Education}

Morten Meldal was born in 1954 in Copenhagen, Denmark.

\subsection*{Career and Major Research}

Meldal received his doctorate in chemistry from the University of Copenhagen in 1982. He worked at the Carlsberg Laboratory in Copenhagen and became a professor at the University of Copenhagen. In 2002, working with solid-phase peptide chemistry, he developed the copper-catalyzed azide-alkyne cycloaddition (CuAAC), which joins an azide and an alkyne rapidly and selectively to form a triazole ring.

\subsection*{Nobel Prize-winning Work}

The work for which Morten Meldal was awarded the Nobel Prize in Chemistry in 2022 was described by the Nobel Committee as: "for the development of click chemistry and bioorthogonal chemistry".

Meldal's CuAAC reaction, developed on solid support, turned the azide-alkyne cycloaddition into a fast, reliable, and water-compatible reaction, the archetypal click reaction.

\subsection*{Significance and Impact}

Because the reaction works quickly, cleanly, and with nearly any partners, it is used to join molecules in combinatorial chemistry, drug discovery, and materials science.

\subsection*{Later Career and Honors}

Meldal continues his research in chemistry and chemical biology at the University of Copenhagen.

\subsection*{Legacy}

CuAAC is one of the most widely used reactions in modern chemistry, applied from the modification of biomolecules to the synthesis of new materials.

\subsection*{Modern Developments}

Meldal's independent development of copper-catalyzed azide-alkyne cycloaddition, the reaction at the heart of click chemistry, provided a robust tool for joining molecules together. Click chemistry is now used in drug discovery, materials science, bioconjugation, and the synthesis of large molecular libraries.

\subsection*{Selected Publications}

"Peptidotriazoles on Solid Phase: [1,2,3]-Triazoles by Regiospecific Copper(I)-Catalyzed 1,3-Dipolar Cycloaddition of Terminal Alkynes to Azides" (Journal of Organic Chemistry, 2002)

\clearpage

\section{K. Barry Sharpless}

\subsection*{Early Life and Education}

K. Barry Sharpless was born in 1941 in Philadelphia, Pennsylvania, USA.

\subsection*{Career and Major Research}

Sharpless earned a bachelor's degree from Dartmouth College in 1963 and a doctorate from Stanford University in 1968. He became a professor at the Scripps Research Institute in La Jolla. He coined the term click chemistry for modular reactions that join molecular building blocks quickly and reliably, and with his colleagues he developed the copper-catalyzed azide-alkyne cycloaddition in 2002. He had earlier received the 2001 Nobel Prize in Chemistry for his work on asymmetric oxidation.

\subsection*{Nobel Prize-winning Work}

The work for which K. Barry Sharpless was awarded the Nobel Prize in Chemistry in 2022 was described by the Nobel Committee as: "for the development of click chemistry and bioorthogonal chemistry".

Sharpless defined click chemistry as reactions that are fast, selective, and high-yielding under simple conditions, and the azide-alkyne reaction he helped develop became its most famous example.

\subsection*{Significance and Impact}

Click chemistry gave chemists a simple way to connect molecules, and it is now used across drug discovery, materials, and biology. The 2022 award made Sharpless a two-time Nobel recipient.

\subsection*{Later Career and Honors}

Sharpless continues his research on click chemistry at the Scripps Research Institute.

\subsection*{Legacy}

Click chemistry and the CuAAC reaction are foundational to modern chemical biology and have been used to build libraries of drug candidates and new materials.

\subsection*{Modern Developments}

Sharpless's formulation of click chemistry, and his independent discovery of the copper-catalyzed azide-alkyne cycloaddition with Meldal, provided chemistry with a general method for rapidly and reliably joining molecular fragments. Click chemistry is now ubiquitous in drug development, chemical biology, and materials science, and his earlier asymmetric epoxidation remains a landmark of enantioselective synthesis.

\subsection*{Selected Publications}

"The First Practical Method for Asymmetric Epoxidation" (Journal of the American Chemical Society, 1980)

\clearpage

\section*{Chapter Summary}

The 2022 prize recognized the development of click chemistry and bioorthogonal chemistry: Sharpless conceived click chemistry and helped develop the copper-catalyzed azide-alkyne reaction, Meldal developed the reaction on solid phase, and Bertozzi made such reactions work inside living organisms. These methods allow molecules to be joined quickly and specifically and are now used in drug discovery and cell biology.
